\documentclass[oneside,a4paper,14pt]{extarticle} %размер шрифта 14
\usepackage[T1,T2A]{fontenc}
\usepackage[
	a4paper,
	letterpaper,
	top=2cm,
	bottom=2cm,
	left=2.5cm,
	right=1.5cm
]{geometry}
\usepackage[utf8]{inputenc} %кодировка текста
\usepackage[russian]{babel} %поддержка русского языка
\usepackage{textcomp} %текстовые символы
\usepackage{indentfirst} %корректировка отступов
\usepackage{listings} %для кода !!
\usepackage{float} % для "впихивания" изображений в странницу
\usepackage{xcolor}
\usepackage{enumitem} %для нумерации буквами
\usepackage{graphicx} %работа с изображениям
\usepackage{mwe} % for blindtext and example-image-a in example
\usepackage{wrapfig}
\usepackage{caption}
\usepackage{amsmath}  % для формул и символов
\usepackage{amsfonts}
\usepackage{amsthm}
\usepackage{graphicx}
\usepackage[all]{xy}
\usepackage[breaklinks]{hyperref}
\usepackage{xurl}

%%размеркегляузаголоковразделов
\usepackage{titlesec}
\titleformat{\section}
{\normalsize\bfseries}
{\thesection} {1em}{}
\titleformat{\subsection}
{\normalsize\bfseries}
{\thesubsection} {1em}{}
\titleformat{\subsubsection}
{\normalsize\bfseries}
{\thesubsection} {1em}{}

\renewcommand\baselinestretch{1.33}\normalsize % межстрочный интервал
\setlength{\parindent}{1.25cm}
\usepackage{indentfirst}

\lstset{
  language=C,
  basicstyle=\ttfamily\small,
  keywordstyle=\color{blue},
  commentstyle=\color{green!60!black},
  stringstyle=\color{red},
  numbers=left,
  numberstyle=\tiny\color{gray},
  frame=single,
  breaklines=true,
  showstringspaces=false,
  tabsize=2,
  inputencoding=utf8,
  extendedchars=true
}

\begin{document}
	\newpage\thispagestyle{empty}
	\begin{center}
		МИНИСТЕРСТВО НАУКИ И ВЫСШЕГО ОБРАЗОВАНИЯ\\
			РОССИЙСКОЙ ФЕДЕРАЦИИ
			ФЕДЕРАЛЬНОЕ ГОСУДАРСТВЕННОЕ БЮДЖЕТНОЕ\\
			ОБРАЗОВАТЕЛЬНОЕ
			УЧРЕЖДЕНИЕ ВЫСШЕГО ОБРАЗОВАНИЯ\\
			«ВЯТСКИЙ ГОСУДАРСТВЕННЫЙ УНИВЕРСИТЕТ»\\
			Институт математики и информационных систем\\
			Факультет автоматики и вычислительной техники\\
			Кафедра электронных вычислительных машин
	\end{center}
	\vspace{20mm}

	\begin{center}
		Отчёт по лабораторной работе №6\\
		по дисциплине\\
		<<Алгоритмы и структуры данных>>\\
	\end{center}
	\vspace{48mm}

	\noindent
	Выполнил студент гр. ИВТб-1303-06-00 \hfill \rule[-0,5mm]{30mm}{0.15mm}\,/Гортоломей И.К./
	\newline
	Проверил заведующий кафедрой ЭВМ \hfill \rule[-0,5mm]{30mm}{0.15mm}\,/Долженкова М.Л./

	\vfill
	\begin{center}
		Киров\\
		2026
	\end{center}

	\newpage
	% \thispagestyle{empty}

	\section*{Цель работы:}
	Отточить навыки решения задач в условиях ограниченного времени.

  \section*{Вводные}
  Для выполнения данной лабораторной работы был предоставлен \\ Яндекс.Контест
  \url{https://official.contest.yandex.ru/contest/94203}

  \section*{Описание процесса решения}
  Был запущен контест, который длится 2 часа.
  Я приступил к решению первой задачи: A. Запрещенные биграммамы
  (\url{https://official.contest.yandex.ru/contest/94203/problems/A/}).
  На сколько я понял, в задаче требуется найти кол-во вариантов строк N длинны,
  чтобы в них небыло повторений одного и того же символа два раза подряд.
  То есть, у нас просто чередуются множители 3 и 1, как: 1x3x1x3x1... если начинается с гласной
  и 3x1x3x1x3... в случае если N - чётное, то резульат будет $2 \cdot 3^{n\ mod\ 2}$
  иначе в случае нечётного N будет  $4 \cdot 3^{n\ mod\ 2}$. Также результат нужно вывести в
  виде остатка от деления на $10^9 + 7$. \\

  \lstinputlisting[language=Python, caption={Решение задачи на Python}]{1.py}
  Решение данной задачи заняло примерно 30 минут.

  \newpage
  
  Далее я приступил к решению задачи B. RL-агент стремится к хаосу
  (\url{https://official.contest.yandex.ru/contest/94203/problems/B/}). 
  В данной задаче нужно определить, есть ли в массиве разница между i-ым элементом и индексом i
  больше k, если да то вывести 1 иначе 0. Чтобы не проверять весь массив после каждой  перестановки
  я решил проверять каждую перестановку, сделала ли перестановка массив <<хаотичным>>. \\

  \lstinputlisting[caption={Решение задачи на C}]{2.c}
  Решение данной задачи заняло тоже примерно 30 минут. \\

  После решения первых двух задач остался 1 час соревнования. Я прочитал задание
  C. Рекомендательная система доставки
  (\url{https://official.contest.yandex.ru/contest/94203/problems/C/}).
  По скольку у нас артикулов 100 тыс. то можно заранее выделить под них память, что я и сделал.
  Массив из 100 тыс. элементов хранит число int - стоимость доставки и число int - номер склада.
  Стоимость доставки установил в максимум, а номер склада в 0. (так как нумерация начинается с 1).
  Далее при вводе складов, проверяем каждый артикул по массиву: если стоимость доставки явно меньше
  чем уже записанная в массиве или если стоимость такая же, но номер склада
  меньше - то заменяем запись. Потом переходим к обработке К запросов. Каждый запрос содержит артикул
  и стоимость товара. если по артикулу в массиве у нас номер склада равен 0 - то такого товара нет,
  выводим -1 -1. Иначе выводим стоимость + цену доставки из массива.
  \newpage
  \lstinputlisting[caption={Решение задачи на C}]{3.c}
  Решение третьей задачи заняло тоже примерно 30 минут. \\

  Далее я открыл задачу D. Метрика согласованности кластеризаторов
  (\url{https://official.contest.yandex.ru/contest/94203/problems/D/}).
  Я прочитал задачу несколько раз что заняло 15 минут, но не понял как её решать. Текст сильно
  запутанный, а формулы обфусцированы. Отслеживание всех логических цепочек для понимая поставленной
  задачи слишком сложно при напряженни из-за ограниченного времени. \\

  В оставшиеся 15 минут я решил посмотреть задачи: <<E. Ты узнаешь ее из тысячи>>
  и <<F. Автостопом по Галактике>>. я решил что не успею их решить за 15 минут. Я считаю, что эти
  задачи вообще не уместные для контеста, если я правильно понял что в них нужно обучать модель.

  \section*{Выводы}
  Подводя итоги можно сказать, что для успешного решения задач в контесте нужно уметь быстро 
  анализировать саму суть задачи, отбрасывая лишнее описание задачи. Как показала практика, не
  всегда удаётся быстро понять задачу, как например в задании D. До этой задачи успешно получалось,
  потому что в задаче описание было простое и был приведён пример, по которому можно быстро понять,
  но в случае задачи D - пример не однозначный, а описание задачи слишком сложное.
\end{document}
